\scanpage{35}\mspage{28}
% archive sheet number: 34
\uline{Remarques} a) Une structure de Quillen exacte à gauche ou exacte à droite est exacte.

b) Si on a une \deleted{structure de} donnée $(C,TC,F,TF)$ satisfaisant $\mathrm M1,\mathrm M2,\mathrm M',\mathrm M''$, quand on sature $C,TC$ à gauche et $F,TF$ à droite, pour les stabilités précisées dans $\unsure{\mathrm M4|\mathrm M1},\mathrm M3$, on trouve une structure de Quillen. (Tous les axiomes sont évidents, sauf $\mathrm M'$, pour lequel je renvoie à p. 9, 10, \unsure{notamment} prop. 4.) Si on avait $\mathrm M2$ bis, \unsure{cel|cela} sera encore vérifié (a fortiori) après saturation, i.e. on a une structure de Quillen exacte.

\uline{Définition 1 bis} La structure de Quillen faible est dite \uline{structure de Quillen} (tout court) si elle est exacte à gauche et à droite, ou (comme on dit) ``biexacte''\unsure{;} \inserted{alors $(C,F)$ déterminent déjà $(C,TC,F,TF)$} (Ça correspond exactement aux ``model categories'' de Quillen.) On dit que la structure de Quillen faible est \uline{close} si on a
\[
\begin{aligned}
&C\mathrel{\leftrightarrow\mkern-15mu\mid\mkern4mu}TF,\qquad TC\mathrel{\leftrightarrow\mkern-15mu\mid\mkern4mu}F,\\
&\deleted{\unsure{TC}}TF=C_{*},\ C=(TF)^{*}\qquad\text{et}\quad F=(TC)_{*},\ TC=F^{*}.
\end{aligned}
\]
% slanted
\margin{Cela signifie aussi que $C,TF,TC,F$ sont stables par facteurs directs. On dit que la str. $\unsure{Qf|Q}$ est close à g. (resp. à dr.) si $C,TC$ (resp. $F,TF$) \unsure{sont} clos à g. (resp. à dr.)}

Une structure de Quillen faible à la fois biexacte (i.e. pas faible du tout !) et close est appelée str. de Qu. \uline{parfaite}. (Ça correspond exactement \unsure{à} \deleted{\unsure{aux|au}} ``closed model category'' de Quillen)
